% 第5章 5.1.2：五点 x=-1,-0.5,0,0.5,1 上的 Gram 矩阵热图与特征值。
% (a) 高斯核 sigma=0.5；(b) 仿射核 1+xz；(c) 两者的特征值。
% 数值由同目录 generate.py 计算：高斯特征值 2.2171,1.5314,0.8199,0.3329,0.0987；仿射 5,2.5,0,0,0。
\begin{tikzpicture}[font=\figurefont\footnotesize]
\def\cell{5.6mm}
% (a) Gaussian heatmap
\begin{scope}
\node[anchor=south] at (2.5*\cell,5*\cell+1.2mm) {(a) 高斯核};
\foreach \row [count=\i from 0] in {
  {1.0000,0.6065,0.1353,0.0111,0.0003},
  {0.6065,1.0000,0.6065,0.1353,0.0111},
  {0.1353,0.6065,1.0000,0.6065,0.1353},
  {0.0111,0.1353,0.6065,1.0000,0.6065},
  {0.0003,0.0111,0.1353,0.6065,1.0000}} {
  \foreach \v [count=\j from 0] in \row {
    \pgfmathsetmacro{\p}{100*\v}
    \fill[FigureBlue!\p!white] (\j*\cell,{(4-\i)*\cell}) rectangle ++(\cell,\cell);
    \pgfmathsetmacro{\shown}{\v<0.005 ? 0 : \v}
    \pgfmathparse{\v>0.55?1:0}\ifnum\pgfmathresult=1 \def\tcol{white}\else\def\tcol{FigureInk}\fi
    \node[font=\figurefont\tiny,text=\tcol] at ({(\j+0.5)*\cell},{(4.5-\i)*\cell}) {\pgfmathprintnumber[fixed,precision=2,zerofill]{\shown}};
  }
}
\draw[FigureStroke,line width=.5pt] (0,0) rectangle (5*\cell,5*\cell);
\end{scope}
% (b) affine heatmap
\begin{scope}[xshift=5*\cell+7mm]
\node[anchor=south] at (2.5*\cell,5*\cell+1.2mm) {(b) 仿射核};
\foreach \row [count=\i from 0] in {
  {2.00,1.50,1.00,0.50,0.00},
  {1.50,1.25,1.00,0.75,0.50},
  {1.00,1.00,1.00,1.00,1.00},
  {0.50,0.75,1.00,1.25,1.50},
  {0.00,0.50,1.00,1.50,2.00}} {
  \foreach \v [count=\j from 0] in \row {
    \pgfmathsetmacro{\p}{50*\v}
    \fill[FigureBlue!\p!white] (\j*\cell,{(4-\i)*\cell}) rectangle ++(\cell,\cell);
    \pgfmathparse{\v>1.1?1:0}\ifnum\pgfmathresult=1 \def\tcol{white}\else\def\tcol{FigureInk}\fi
    \node[font=\figurefont\tiny,text=\tcol] at ({(\j+0.5)*\cell},{(4.5-\i)*\cell}) {\pgfmathprintnumber[fixed,precision=2,zerofill]{\v}};
  }
}
\draw[FigureStroke,line width=.5pt] (0,0) rectangle (5*\cell,5*\cell);
\end{scope}
% (c) eigenvalues
\begin{axis}[at={(10*\cell+17mm,0)},anchor=south west,width=46mm,height=36mm,
  ybar,bar width=1.6mm,area legend,axis lines=left,ymin=0,ymax=5.4,xmin=0.4,xmax=5.6,
  xtick={1,2,3,4,5},ytick={0,1,2,3,4,5},xlabel={序号$k$},xlabel style={yshift=1.5mm},
  x axis line style={draw=FigureStroke,line width=.6pt},y axis line style={draw=FigureStroke,line width=.6pt},
  tick style={draw=FigureStroke},tick label style={font=\figurefont\scriptsize,text=FigureInk},
  label style={font=\figurefont\scriptsize},
  ymajorgrids,grid style={draw=FigureGrid,line width=.4pt},
  title={(c) 特征值},title style={font=\figurefont\footnotesize},
  legend style={draw=none,fill=none,font=\figurefont\scriptsize,at={(0.98,0.98)},anchor=north east,cells={anchor=west}}]
\addplot[fill=FigureBlue,draw=FigureBlue] coordinates {(1,2.2171) (2,1.5314) (3,0.8199) (4,0.3329) (5,0.0987)};
\addplot[fill=FigureYellow,draw=FigureYellow] coordinates {(1,5) (2,2.5) (3,0) (4,0) (5,0)};
\legend{高斯核,仿射核}
\end{axis}
\end{tikzpicture}
